\documentclass[10pt,a4paper]{amsart}
\usepackage[margin=19mm]{geometry}
\usepackage{amsmath,amssymb,mathtools}
\usepackage[scr=rsfs,cal=euler]{mathalfa}
\usepackage{xfrac}
\usepackage[unicode,hidelinks,bookmarks=false]{hyperref}
\newcommand{\ie}{i.e.\ }
\newcommand{\cf}{cf.\ }
\newcommand{\resp}{respectively\ }
\newcommand{\Subsection}{\subsection}
\providecommand{\nobreakdash}{\nobreak\hbox{-}}
\setlength{\emergencystretch}{2em}
\multlinegap0pt
\usepackage{amssymb}
\usepackage{mathtools}
\usepackage{xcolor}
\newtheorem{lemma}{Lemma}
\title{Syntomic formalism with coefficients}
\author{Fabrizio Andreatta, Massimo Bertolini, Marco Adamo Seveso, Rodolfo Venerucci}
\date{Source: arXiv:2512.02905v2 (CC BY 4.0)}
\begin{document}
\maketitle

\noindent\emph{Source and licence.} Syntomic formalism with coefficients. Version arXiv:2512.02905v2 (CC BY 4.0), distributed by arXiv under the Creative Commons Attribution 4.0 International licence, http://creativecommons.org/licenses/by/4.0/, as recorded in the arXiv licence metadata for this exact version and shown on the abstract page https://arxiv.org/abs/2512.02905v2. Full licence notice: \texttt{licenses/arxiv-2512.02905-CC-BY-4.0.txt}.\par\smallskip

\noindent\textit{Editorial scope.} Counted unit: the article's lemma L1 (source0.tex lines 1379--1392). The article's complete original proof of it is delivered with it (source0.tex lines 1394--1470, one \texttt{proof} environment of 77 lines). No mathematics is written, repaired or re-derived: the statement and the proof are the article's own lines, in the article's own notation, and no retained line is substituted or re-flowed.  No further source-local context is needed: every cross-reference of every retained line resolves inside the two delivered spans.  Nothing was omitted from any retained span, so the package carries no omission record.  No retained line contains a citation, so the wrapper carries no bibliography and no bibliography entry.  No retained line contains a figure environment, an image-inclusion command or drawing code, so no image is delivered and the package declares no asset dependency.  Editorial formatting only, in addition: an AMS \texttt{amsart} preamble that restates the article's own declarations and macros used by the retained lines, namely the environments \texttt{lemma} on separate counters, together with the standard packages those lines need.  The retained environments print in this document as Lemma 1 in order of appearance, whereas the article prints the same items with the numbering of its own full document.  The complete native source of the article as distributed in the arXiv e-print for this exact version is bundled unchanged as \texttt{sources/190224/source0.tex}.
\begin{lemma}\label{lemma:Kos L1} 
Let $d_{1},\ldots d_{n}\in \mathrm{End}_{R}(
A) $  be endomorphisms such that for $j=1,\ldots,n_{o}$ we have  $d_{i}\left( x_{j}\right) =0$ if $i\neq j$ and $d_{j}\left( x_{j}\right) =x_{j}$. Let $ \nabla _{1},\ldots,\nabla _{n}\in \mathrm{End}_{R}(M) $ be commuting homomorphisms such that, for every $i=1,\ldots,n$,we have
\begin{equation}
\nabla _{i}\left( am\right) =d_{i}\left( a\right) m+a\nabla _{i}\left(
m\right) \text{ for every }a\in A\text{ and }m\in M  \label{Kos L1 F1 Clm}
\end{equation}

Assume that the residue $\mathrm{Res}(\nabla _{i})\in \mathrm{End}_R(M/x_i M)$ is
nilpotent for every $i=1,\ldots,n_{o}$.  Then  $I_{S}K^\ast\left( 
\widehat{\nabla }_{1},\ldots,\widehat{\nabla }_{n}\right) $ is an acyclic
subcomplex of $K^\ast\left( \widehat{\nabla }_{1},\ldots,\widehat{
\nabla }_{n}\right) $.  
\end{lemma}

\begin{proof} We remark that due our assumption $\nabla_i(x_i M)\subset x_i M$ so that $\mathrm{Res}(\nabla _{i})$ is well defined. Using (\ref{Kos L1 F1 Clm}) and the assumption that $ d_{i}\left( x_{j}\right) =0$ or $x_{j}$ if $i=j$ for $j\in \left\{ 1,\ldots,n_{o}\right\} $, we deduce from Property (v) that $I_{S}K^\ast\left( \widehat{\nabla }_{1},\ldots,\widehat{\nabla }
_{n}\right) \subset K^\ast\left( \widehat{\nabla }_{1},\ldots,\widehat{
\nabla }_{n}\right) $ is a subcomplex.  


Notice that  $K^\ast\left( \widehat{\nabla }_{1},\ldots,\widehat{\nabla }
_{n}\right) =K^\ast\left( \nabla _{1},\ldots,\nabla _{n}\right) $ when $M$
is $I$-adically complete and separated by Prperty (vi). Since we have the equality of $I$-adic completions $\widehat{\frac{M}{I_{\left\{
i\right\} }M}}\simeq \frac{\widehat{M}}{I_{\left\{ i\right\} } \widehat{M}}$
the nilpotency of $\mathrm{Res}\left( \nabla _{i}\right) $
implies the nilpotency of $\mathrm{Res}\left( \widehat{\nabla }_{i}\right) $. 
In order to prove the claim we can then assume that $M$ is already $I$
-adically complete and separated. Setting $r:=\left\vert S\right\vert $, we
will prove our claim by a double induction.

\emph{Step 1: }$\left( r,n_{o},n\right) =\left( 1,1,1\right) $. The generic
element $z$ of $x_{1}K^{0}\left( \nabla _{1}\right) $ is of the form $
z=x_{1}z_{1}$ where $z_{1}\in K^{0}\left( \nabla _{1}\right) $\ and, setting 
$N:=\nabla _{1}-{\rm Id}$, we have $N\left( z\right) =x_{1}\nabla _{1}\left( z_{1}\right) $. By
induction, we find that $N^{t}\left( z\right) =x_{1}\nabla _{1}^{t}\left(
z_{1}\right) $ for every $t\in \mathbb{N}_{\geq 1}$. Our assumption that $
\mathrm{Res}\left( \nabla _{1}\right) $ is nilpotent implies that there
exists some $t\in \mathbb{N}_{\geq 1}$ such that $\nabla _{1}^{t}\left(
y\right) \in x_{1}M$ for every $y\in M$. Hence, $N^{t}\left(
z\right) \in x_{1}^{2}M$ for every $z\in x_{1}K^{0}\left( \nabla _{1}\right) 
$ or, equivalently, considering $End_{R}\left( x_{1}M\right) $ as a left $A$
-module via the composition with the multiplication by elements of $A$ on
the codomain, then $N^{t}\in
x_{1}End_{R}\left( x_{1}M\right) $, implying that $N$ is topologically
nilpotent in the $I$-adic topology. By a geometric series argument, we deduce that
\begin{equation*}
\nabla _{1}=1+ N_1\in A^{\times }\left( 1+ x_1 \mathrm{End}_{R}\left(
x_{1}M\right) \right) \subset \mathrm{Aut}_{R}\left( x_{1}M\right) \text{.}
\end{equation*}
is an isomorphism, proving the claim in this case.

\emph{Step 2: }$\left( r,n_{o},n\right) =\left( 1,n_{o},n\right) $. We may
assume without loss of generality that $S=\left\{ 1\right\} $ and we prove
our assertion by induction on $n$, the case $n=1$ being discussed in Step 1. 
Suppose that $n>1$. According to $\left( v\right) $
above, $x_{1}K^\ast\left( \nabla _{1},\ldots,\nabla _{m}\right) =K^{\ast
}\left( \nabla _{1\mid x_{1}M},\ldots,\nabla _{m\mid x_{1}M}\right) $ for $m\in
\left\{ n-1,n\right\} $ and, hence, thanks to Property (iii) we know
that $x_{1}K^\ast\left( \nabla _{1\mid
x_{1}M},\ldots,\nabla _{n\mid x_{1}M}\right) $ is the cone of
$$ -\nabla
_{n\mid x_{1}M}:x_{1}K^\ast\left( \nabla _{1\mid x_{1}M},\ldots,\nabla
_{n-1\mid x_{1}M}\right) \rightarrow x_{1}K^\ast\left( \nabla _{1\mid
x_{1}M},\ldots,\nabla _{n-1\mid x_{1}M}\right)$$twisted by $-1$. Using the inductive step one deduces that  $
x_{1}K^\ast\left( \nabla _{1\mid x_{1}M},\ldots,\nabla _{n-1\mid
x_{1}M}\right) $ is acyclic and we conclude the claimed acyclicity of $
x_{1}K^\ast\left( \nabla _{1\mid x_{1}M},\ldots,\nabla _{n\mid
x_{1}M}\right) $ using the identification with the cone.

\emph{Step 3: arbitrary }$\left( r,n_{o},n\right) $\emph{\ with }$1\leq
r\leq n_{o}\leq n$. We may assume without loss of generality that $S=\left\{
1,\ldots,r\right\} $ and we prove our assertion by induction on $r$, the case $
r=1$ being discussed in Step 2. Suppose that $r>1$ and consider
the exact sequence
\begin{equation*}
0\longrightarrow I_{\left\{ 1,\ldots,r-1\right\} }M\longrightarrow I_{\left\{
1,\ldots,r\right\} }M\longrightarrow \frac{I_{\left\{ 1,\ldots,r\right\} }M}{
I_{\left\{ 1,\ldots,r-1\right\} }M}\longrightarrow 0\text{.}
\end{equation*}
Set $M^{\prime \prime }:=\frac{M}{I_{\left\{ 1,\ldots,r-1\right\} }M}$. Then $
x_{r}M^{\prime \prime }=\frac{I_{\left\{ 1,\ldots,r\right\} }M}{I_{\left\{
1,\ldots,r-1\right\} }M}$. Write $\nabla
_{i}^{\prime \prime }$ for the operator induced by $\nabla_i$ on $x_{r}M^{\prime \prime }$.
Thanks to Property (iv)  we have an exact sequence of complexes of $R$-modules expressing $K^\ast\left( \nabla _{1\mid I_{\left\{
1,\ldots,r\right\} }M},\ldots,\nabla _{n\mid I_{\left\{ 1,\ldots,r\right\} }M}\right)$ as an extension of 
$K^\ast\left( \nabla _{1}^{\prime \prime },\ldots,\nabla
_{n}^{\prime \prime }\right)$ by $K^\ast\left( \nabla _{1\mid I_{\left\{
1,\ldots,r-1\right\} }M},\ldots,\nabla _{n\mid I_{\left\{ 1,\ldots,r-1\right\}
}M}\right) $.  By induction, we know that the latter is acyclic. Hence, it suffices to show that  that $K^\ast\left( \nabla _{1}^{\prime \prime },\ldots,\nabla
_{n}^{\prime \prime }\right)$ is acyclic. This follows from the case $\left(
r,n_{o},n\right) =\left( 1,n_{o},n\right) $ discussed in Step 2.
\end{proof}

\end{document}
